\task{} aims to infer machinery operation states from time-indexed sensor records, supporting operation monitoring, efficiency evaluation, and data-driven farm management.
GNSS trajectory methods are effective for field-road discrimination and coarse operational recognition, but they remain limited for fine-grained combine-harvester operations that share overlapping motion patterns in position, speed, heading, curvature, and dwell time.
For these ambiguous states, discriminative evidence often lies in complementary sensor streams: audio captures engine load, idling, unloading, and working-component engagement, whereas images capture crop contact, material flow, harvesting scenes, and field-road context.
However, multimodal use is nontrivial because audio lacks path geometry, images are sensitive to illumination, viewpoint, and occlusion, and the three streams differ in sampling rate, noise characteristics, temporal alignment, and feature space.
This paper studies an eleven-class combine-harvester setting and proposes \ours{}, a multimodal model that integrates trajectory, ViT visual, and audio streams for \task{}.
\ours{} encodes causal GNSS trajectories, synchronized ViT visual frames, and acoustic spectrograms with modality-specific branches, then applies class-adaptive logit-level fusion to combine motion, scene, and process cues under a unified classification objective.
Experiments are conducted on a real-world dataset collected over three months of maize harvest operations in Ningxia, Shaanxi, and Inner Mongolia.
On the held-out test set, \ours{} achieves 81.36\% accuracy, 63.64\% macro-F1, and 81.01\% weighted-F1, exceeding the strongest unimodal baseline by 9.00 macro-F1 points.
Compared with direct trimodal concatenation, the class-adaptive gate improves accuracy, macro-F1, and weighted-F1 by 2.21, 2.44, and 2.56 percentage points, respectively.
Per-class analysis shows that \ours{} improves 8 of 11 operation classes over the best single modality, with large gains for reverse transfer, turning transfer, and idling waiting, while unloading and road driving remain challenging.
To the best of our knowledge, this is the first synchronized trajectory, ViT visual, and audio study of eleven-class \task{}.
